\documentclass[11pt]{article}
\usepackage[a4paper,margin=26mm]{geometry}
\usepackage{amsmath,amssymb,amsthm,mathtools}
\usepackage[hidelinks]{hyperref}
\newtheorem{theorem}{Theorem}
\newtheorem{lemma}[theorem]{Lemma}
\newcommand{\ML}{E}
\newcommand{\dd}{\,\mathrm{d}}
\title{Complete monotonicity of a squared Mittag--Leffler modulus}
\author{Research manuscript}
\date{October 28, 2026}
\begin{document}
\maketitle
\begin{abstract}
We prove complete monotonicity of
$|E_{2a,1}(-t)|^2+t|E_{2a,1+a}(-t)|^2$ on $(0,\infty)$ for
$0<a\leq1/2$, the assertion of Conjecture~3 in the preprint of
Chorfi--Et-tahri--Maniar--Yamamoto. The proof identifies a positive
one-dimensional Laplace representation for the half-order cross kernel,
then transfers it to arbitrary $a$ through finite normalized Beta averages
and Euler's limit for the Gamma function. Classical integral and convexity
tools are used; no new general closure principle is claimed.
\end{abstract}

\section{Statement and the half-order kernel}
For $a,b>0$, let $E_{a,b}(z)=\sum_{n\geq0}z^n/\Gamma(an+b)$.
\begin{theorem}\label{thm:main}
For $0<a\leq1/2$, the function
\[
 f_a(t)=|E_{2a,1}(-t)|^2+t|E_{2a,1+a}(-t)|^2
\]
is completely monotone: it is smooth on $(0,\infty)$ and
$(-1)^k f_a^{(k)}(t)\geq0$ for every integer $k\geq0$ and $t>0$.
\end{theorem}
The theorem is precisely Conjecture~3, equation~(5.2), in
\cite[Section~5]{CEMY}. The other zero and evolution-equation questions
in that source are not addressed here. The genuine even/odd series identity
\[
 E_{a,1}(z)=E_{2a,1}(z^2)+zE_{2a,1+a}(z^2)
\]
and the reality of the two series at negative real arguments give
$f_a(t)=|E_{a,1}(i\sqrt t)|^2$ for $t\geq0$.

Put $I(x)=\int_0^1e^{xu^2}\dd u$. Termwise integration of the exponential
series, dominated on $[0,1]$ by $e^{|x|}$, gives
\begin{equation}\label{eq:half}
 \int_0^1e^{-x(1-u^2)}\dd u
 =\frac{\sqrt\pi}{2}E_{1,3/2}(-x),\qquad
 E_{1/2,1}(ir)=e^{-r^2}\left(1+\frac{2ir}{\sqrt\pi}I(r^2)\right).
\end{equation}
For completeness, the polynomial moments used in the first identity are
\[
 B_n=\int_0^1(1-u^2)^n\dd u
 =\frac{\sqrt\pi\,n!}{2\Gamma(n+3/2)}.
\]
Integrating the derivative of $u(1-u^2)^{n+1}$ yields
$B_0=1$ and $B_{n+1}=2(n+1)B_n/(2n+3)$; the displayed formula follows
from the Gamma recurrence. The second identity in \eqref{eq:half} follows
from the even/odd decomposition and $E_{1,1}(z)=e^z$.

For $v,w\geq0$ and $t\geq0$, define
\[
 K(v,w,t)=\operatorname{Re}\!\left[
 E_{1/2,1}(i\sqrt{vt})\overline{E_{1/2,1}(i\sqrt{wt})}\right].
\]
Writing $S=v+w$, equation~\eqref{eq:half} gives
\begin{equation}\label{eq:product}
 K(v,w,t)=e^{-St}\left[1+\frac{4t\sqrt{vw}}{\pi}I(vt)I(wt)\right].
\end{equation}
\begin{lemma}\label{lem:kernel}
For $v,w>0$, the expression in \eqref{eq:product} equals
\begin{equation}\label{eq:laplace}
 \frac{2\sqrt{vw}}{\pi}\left[
 \int_0^1\frac{e^{-tw(1-u^2)}}{v+wu^2}\dd u+
 \int_0^1\frac{e^{-tv(1-u^2)}}{w+vu^2}\dd u\right].
\end{equation}
In particular, $K(v,w,\cdot)$ is a positive Laplace mixture.
For a zero scale, $K(v,w,t)=e^{-(v+w)t}$.
\end{lemma}
\begin{proof}
Denote the right side of \eqref{eq:laplace} by $J(t)$. Differentiation of
$u e^{xu^2}$ and integration give $I(x)+2xI'(x)=e^x$. The product rule
therefore shows that the expression $P(t)$ in \eqref{eq:product} satisfies
\[
 P'(t)+SP(t)=\frac{2\sqrt{vw}}{\pi}
 \left[e^{-wt}I(wt)+e^{-vt}I(vt)\right].
\]
The same equation holds for $J$: differentiation under its compact
integrals cancels each denominator, since
$S-w(1-u^2)=v+wu^2$ and $S-v(1-u^2)=w+vu^2$.
All integrands and their parameter derivatives are continuous and bounded
on compact rectangles. Finally,
\[
 J(0)=\frac2\pi\left[\arctan\sqrt{w/v}+\arctan\sqrt{v/w}\right]=1=P(0).
\]
Thus $e^{St}(P-J)$ has zero derivative and vanishes at zero.
This proves the identity. Both weights and both rates in
\eqref{eq:laplace} are nonnegative. The zero-scale formula follows
directly from \eqref{eq:product}.
\end{proof}

\section{Finite averaging and arbitrary order}
Fix $\beta=2a\in(0,1]$. For $k\geq1$, use the normalized positive averaging
operator
\[
 (T_k g)(x)=\beta g(x)+(1-\beta)k\int_0^1u^{k-1}g(xu)\dd u.
\]
It preserves constants and positivity, and is a contraction in the uniform
norm on every interval $[0,C]$. For $q\geq0$,
\[
 T_k(x^q)=\frac{k+\beta q}{k+q}x^q.
\]
For $N\geq1$, let $A_N$ be the composition of $T_1,\ldots,T_{N+1}$,
evaluated at $C_N=N^{1-\beta}$. Then
\begin{equation}\label{eq:moments}
 R_N(q):=A_N(x^q)=N^{(1-\beta)q}
 \prod_{k=1}^{N+1}\frac{k+\beta q}{k+q}
 \longrightarrow\frac{\Gamma(1+q)}{\Gamma(1+\beta q)}.
\end{equation}
This is Euler's Gamma limit. Moreover, $R_N(q)$ increases with $N$.
Indeed, with $d=(1-\beta)q\geq0$,
\[
 \frac{R_{N+1}(q)}{R_N(q)}
 =\left(\frac{N+1}{N}\right)^d\frac{N+2+\beta q}{N+2+q}\geq1.
\]
Here $((N+1)/N)^d\geq1+d/(N+1)$ follows from
$\log((N+1)/N)\geq1/(N+1)$ and $e^y\geq1+y$; multiplying out leaves
$d(1+\beta q)/(N+1)\geq0$.
Consequently $0<R_N(q)\leq\Gamma(1+q)/\Gamma(1+\beta q)$.

Define the genuine complex transforms
\[
 M_N(z)=A_N\!\left(E_{1/2,1}(z\sqrt x)\right)
 =\sum_{n\geq0}\frac{R_N(n/2)}{\Gamma(1+n/2)}z^n.
\]
The equality follows from uniform convergence on the finite interval
$[0,C_N]$. By \eqref{eq:moments},
\[
 0\leq\frac{R_N(n/2)}{\Gamma(1+n/2)}
 \leq\frac1{\Gamma(1+an)},\qquad M_N(z)\longrightarrow E_{a,1}(z).
\]
The entire series on the right supplies a summable majorant on every
bounded complex disk, so these are actual transform limits.

Real linearity, complex conjugation and two independent applications of
$A_N$ show that
\[
 |M_N(i\sqrt t)|^2=A_N^{(v)}A_N^{(w)}K(v,w,t).
\]
This identity is an identity of finite weighted integrals; no existence
of a stable distribution is needed. Lemma~\ref{lem:kernel} makes the
right side a positive average of exponential Laplace kernels. For arbitrary
positive steps $h_1,\ldots,h_k$, every mixed lower difference is nonnegative,
because that of $e^{-\ell t}$ is
\[
 e^{-\ell t}\prod_{j=1}^k(1-e^{-\ell h_j})\geq0\qquad(\ell\geq0).
\]
Finite differences commute with these averages and pass to pointwise
limits. Hence all positive mixed lower differences of $f_a$ are nonnegative.
The function $f_a$ is smooth by its entire real-series expression. Mixed
differences of a smooth function are antitone when one additional lower
difference is nonnegative. Taking their negative derivatives and iterating
gives $(-1)^k f_a^{(k)}\geq0$. This proves Theorem~\ref{thm:main}.

\section{Certificate and research status}
The Lean endpoint is \texttt{MittagModulus.source\_conjecture\_370772}, a
closed proof of the literal two-parameter-series statement. The certificate
proves the series convergence, real and complex analyticity, exact finite
moments, actual averaging and squared-average identities, the kernel
formula including zero scales, and the finite-difference-to-derivative
passage. Its pinned environment is Lean~4.34.1 and Mathlib
\texttt{d13f23b723b}. A fresh isolated source replay passed all 17 modules and 155 declarations,
with warnings as errors and only the standard axioms \texttt{propext},
\texttt{Classical.choice}, and \texttt{Quot.sound}. No independent
formal-review pass is asserted in this draft.

This manuscript and candidate have zero discovery credit pending source,
substantive, prior-art and independent final review, followed by the final
packet replay. The preprint's actual conjecture page has been read; the
published journal body has not yet been compared. The source is also listed
as Chaos Nonlinear Sci. Numer. Simul.\ \textbf{157} (2026), 109758,
doi:10.1016/j.cnsns.2026.109758. No worldwide-priority claim is made.

\begin{thebibliography}{9}
\bibitem{CEMY} S. E. Chorfi, F. Et-tahri, L. Maniar and M. Yamamoto,
\emph{Forward and backward problems for abstract time-fractional
Schr\"odinger equations}, arXiv:2510.03600v1 (2025),
Section~5, Conjecture~3, equation~(5.2).
\href{https://arxiv.org/abs/2510.03600v1}{arxiv.org/abs/2510.03600v1}.
\end{thebibliography}
\end{document}
